% !TEX root = chapter-standalone.tex

\section{Group actions}


While semigroups are like monoids without neutral elements, groups are monoids with additional structure: every element of a group has an inverse.

Monoid actions were defined as monoid morphisms, semigroup actions were defined as semigroup morphisms, so we would expect to find a definition for group actions that follows this pattern. 

To define an action of a group $\grpA$ on a set $\stateset$, an initial idea that one might have is to define it as a group homomorphism of this type
\begin{equation}
\act \colon \grpA \mto \Endof\stateset,
\end{equation}
but the issue we run into is that $\Endof\stateset$ is not a group (in general it contains non-invertible functions). However, the subset $\Autof\stateset \subseteq \Endof\stateset$ of invertible functions on $\stateset$ \emph{is} a group, with function inverses as the group inverses. 
%The notation ``$\Aut$'' makes reference to the word ``automorphism'', which means "invertible morphism" (in this case, invertible functions). 

\begin{ctdefinition}[Group action]
    \label{def:group-cov-action}
    An \maindef{action} of a \SY{group} $\grpA$ on a set $\stateset$ is:

    \constit
    \begin{enumerate}
        \item A \SY{group morphism}
              \begin{equation}
                  \label{eq:act-group}
                  \act \colon \grpA \fto \Autof\stateset.
              \end{equation}
    \end{enumerate}
\end{ctdefinition}

\begin{ctdefinition}[Group action, uncurried]
    \label{def:group-uncurried-action}

    A (right-uncurried) action of a \SY{group} $\grpA = \tup{\grpAset, \mtimesof\grpA, \idgrp_\grpA, \ginv}$ on a set \stateset is:

    \constit
    \begin{enumerate}
        \item a function
              \begin{equation}
                  \label{eq:group-uncurried-action}
                  \act \colon \makecartprod{\grpAset, \stateset} \sto \stateset.
              \end{equation}
    \end{enumerate}
    \condit
    \begin{enumerate}
        \item for all $\statesetel \setin \stateset$ and $\grpelA,\grpelB\setin\grpAset$,
              \begin{equation}
                  \label{eq:group-uncurried-action-1cond}
                  \act(\grpelA \mtimesof\grpA \grpelB, \statesetel) = \act(\grpelB, \act(\grpelA, \statesetel));
              \end{equation}
        \item for all $\statesetel \setin \stateset$,
              \begin{equation}
                  \label{eq:group-uncurried-action-2cond}
                  \act(\idgrp_\grpA, \statesetel) = \statesetel.
              \end{equation}
    \end{enumerate}
\end{ctdefinition}

\begin{remark}
Suppose we have a monoid morphism
\begin{equation}
\act \colon \monoidA \fto \Endof\stateset.
\end{equation}
If $\monoidA$ happens to be a group, then one can prove that the image of $\act$ is necessarily a subset of $\Autof\stateset \subseteq \Endof\stateset$. In other words, any monoid action where the monoid happens to be a group may be viewed as a group action. 
\end{remark}

\todotext{JL: make remark or lemma here that inverses being respected comes for free... and spell out the property in action terms...}



\todotext{\bernina: @JL: add example of vector spaces as actions of a field on an abelian group; mention concept of a module}

\todotext{lemma commented out here}
%\begin{lemma}
%\label{lem:uncurrying-group-actions}
%The data of a \SY{group action}
%\begin{equation}
%\act \colon \grpA \fto \Autof\setA
%\end{equation}
%is equivalent to the data of having a function
%\begin{equation}
%\stylemaps{\alpha} \colon \grpAset \cartprod \setA \sto \setA,
%\end{equation}
%which satisfies the conditions
%\begin{enumerate}
%\item $ \stylemaps{\alpha} (\idgrp_\grpA, \setAel) = \setAel \quad \forall \setAel \setin \setA$;
%\item $ \stylemaps{\alpha} (\grpelb, \stylemaps{\alpha}(\grpela, \setAel)) = \stylemaps{\alpha} (\grpela \mtimes \grpelb, \setAel) \quad \forall \setAel \setin \setA, \forall \grpela, \grpelb \setin \grpAset.
%      $
%\end{enumerate}
%\end{lemma}


\todotext{Add examples, in particular ones relating to group representations and even some definitions perhaps of matrix groups, understood via actions e.g. when notion of symmetry is in the foreground.}



\begin{example}[Time dilation of signals]
In the field of signal-processing, a basic way to model a signal is as a function $f \colon \reals \to \subA$, where the domain models time and the codomain models the possible values that a signal might take. 

In this example, we consider the set of all possible such signals, $\stateset = \subA^\reals$. Let $\tup{\posReals, \cdot, 1}$ be the group of positive reals, with multiplication as composition. We can define a group action $\act \colon \posReals \fto \Autof\stateset$ on signals that models ``time dilation'': for each $\lambda \setin \posReals$,
\begin{equation}
\defmapperiod{\act(\lambda)}{\stateset}{\to}{\stateset}{f}{\bigl(t \mapsto f(t / \lambda)\bigr)}
\end{equation}
\end{example}

\todotextjira{26}{\bernina: @JL: add more exercises}

\todotext{graded exercise commented out here}
%\begin{gradedexercise}[\exname{UncurryingGroupActions}]
%    \label{ex:UncurryingGroupActions}
%    Prove \cref{lem:uncurrying-group-actions}.
%\end{gradedexercise}
%\solutionof{UncurryingGroupActions}

\begin{example}[Invertible matrices acting on vectors]
    \label{exa:matrix-group-action}
    This is the group version of \cref{exa:matrix-action}.
    Fix a natural number $n$, and let $\mgGLn$ denote the set of invertible $\ntimesn$ real matrices.
    Consider the \SY{group} $\grpA = \tup{\mgGLn, \mtimesof\grpA, \idmat, \ginv}$ with composition operation
    \begin{equation}
        \matmor{A} \mtimesof\grpA \matmor{B} \definedas \matmor{B} \matmor{A}
    \end{equation}
    (as in \cref{eq:matrix-action-composition}), neutral element the identity matrix $\idmat$, and inverse operation $\ginv(\matmor{A}) \definedas \matmor{A}^{-1}$.
    This is indeed a \SY{group}: $\matmor{B}\matmor{A}$ is invertible whenever $\matmor{A}$ and $\matmor{B}$ are, associativity and neutrality hold as in \cref{exa:matrix-action}, and $\matmor{A}^{-1}\matmor{A} = \idmat = \matmor{A}\matmor{A}^{-1}$.

    As in \cref{exa:matrix-action}, define, for each $\matmor{A} \setin \mgGLn$,
    \begin{equation}
        \defmapperiod{\act(\matmor{A})}{\reals^n}{\to}{\reals^n}{\vectob{v}}{\matmor{A}\vectob{v}}
    \end{equation}
    Each $\act(\matmor{A})$ is invertible, with inverse $\act(\matmor{A}^{-1})$, so it is an element of $\Autof{\reals^n}$.
    The computation in \cref{exa:matrix-action} shows that $\act(\matmor{A} \mtimesof\grpA \matmor{B}) = \act(\matmor{A}) \mtimesof{\Autof{\reals^n}} \act(\matmor{B})$ for all $\matmor{A}, \matmor{B} \setin \mgGLn$.
    Hence
    \begin{equation}
        \act \colon \grpA \fto \Autof{\reals^n}
    \end{equation}
    is a \SY{group morphism}, that is, an action of $\grpA$ on $\reals^n$.
    In uncurried form, $\act(\matmor{A}, \vectob{v}) = \matmor{A}\vectob{v}$, and condition \cref{eq:group-uncurried-action-1cond} reads $(\matmor{B}\matmor{A})\vectob{v} = \matmor{B}(\matmor{A}\vectob{v})$.
\end{example}


\subsection{Example: ortho-translations}

\begin{definition}
    [General Euclidean group $\mgEn$]
    \label{def:general-euclidean-group}
    The general (real) Euclidean group of order $n$, written $\mgEn$,
    is the \SY{group} given by:

    \constit
    \begin{enumerate}
        \item The set of $(n+1) \times (n+1)$ real matrices of the form
              \begin{equation}
                  \label{eq:En-def}
                  \Ematrix{\matmor{R}}{\vectmor{t}},
              \end{equation}
              where $\matmor{R}\setin \mgOn$ and $\vectmor{t} \setin \reals^n$;
        \item The composition operation given by matrix multiplication;
        \item The neutral element given by the identity matrix;
        \item The inverse operation given by matrix inversion.
    \end{enumerate}
\end{definition}

\begin{definition}
    [Special Euclidean group $\mgSEn$]
    \label{def:special-euclidean-group}
    The special (real) Euclidean group of order $n$, written $\mgSEn$,
    is the \SY{group} given by:

    \constit
    \begin{enumerate}
        \item The set of $(n+1) \times (n+1)$ real matrices of the form
              \begin{equation}
                  \label{eq:SEn-def}
                  \Ematrix{\matmor{R}}{\vectmor{t}},
              \end{equation}
              where $\matmor{R}\setin \mgSOn$ and $\vectmor{t} \setin \reals^n$;
        \item The composition operation given by matrix multiplication;
        \item The neutral element given by the identity matrix;
        \item The inverse operation given by matrix inversion.
    \end{enumerate}
\end{definition}

\begin{table*}
    \caption{Matrix groups}
    \label{tab:matrix-groups}
    \begin{tabular}{cllc}
        \mgGLn & General linear group     & invertible linear transformations \\
        \mgSLn & Special linear group     & invertible linear transformations with determinant equal to 1 \\ \\
        \mgOn  & Orthogonal group         & preserve length of vectors \\
        \mgSOn & Special orthogonal group & rotations \\ \\
        \mgEn  & Euclidean group          & preserve distances and angles                                 & \\
        \mgSEn & Special Euclidean group  & rigid motions \\
    \end{tabular}
\end{table*}

The \SY{groups}~$\mgSEtwo$ and~$\mgSEthree$ are particularly important in robotics because they represent the roto-translations of the plane and 3D space, respectively.

From~\cref{eq:SEn-def} we know we can represent one by a pair~$\tupp{\matmor{R}, \vectmor{t}}$, with~$\matmor{R}\setin\mgSOn$ and~$\vectmor{t} \setin \reals^n$.

If we look at how matrices compose, we get
%
\begin{equation}
    \label{eq:SE-composition}
    \Ematrix{\matmor{R}_2}{\vectmor{t}_2} \Ematrix{\matmor{R}_1}{\vectmor{t}_1} = \Ematrix{\matmor{R}_2 \matmor{R}_1}{\matmor{R}_2\vectmor{t}_1 + \vectmor{t}_2}.
\end{equation}
%
The formula for composition is
%
\begin{equation}
    \label{eq:SE-composition-short}
    \tupp{\matmor{R}_1, \vectmor{t}_1} \mtimesof{\mgSEn} \tupp{\matmor{R}_2, \vectmor{t}_2} = \tupp{\matmor{R}_2 \matmor{R}_1, \matmor{R}_2 \vectmor{t}_1 + \vectmor{t}_2}.
\end{equation}
%

% %
% \begin{equation}
%   \tupp{\matmor{R}_1, \vectmor{t}_1} \mthenof{\mgSEn} \tupp{\matmor{R}_2, \vectmor{t}_2}  = \tupp{\matmor{R}_2 \matmor{R}_1, \matmor{R}_2 \vectmor{t}_1 + \vectmor{t}_2}
% \end{equation}
%
% Note that because we are using $\mthen$ (then), we do not have the nice symmetry
% that we had if we were using $\after$ (after).
% %
% %

The group~$\mgSEn$ induces a transformation on the points of~$\reals^n$.
We are going to call this an \emph{action}.

The action is the following function:
%
\begin{equation}
    \label{eq:apply}
    \defmapperiodset{\mgact}{
        \mgSEn \cartprod \reals^n
    }{
        \reals^n
    }{
        \tupp{ \tupp{\matmor{R}, \vectmor{t}}, \vectob{p}}
    }{
        \matmor{R}\vectob{p} + \vectmor{t}
    }
\end{equation}
Given a roto-translation and a point, the function returns the roto-translated point.
%
We can also see this in matrix form as follows.
We need to substitute for a point~$\vectob{p} \setin \reals^n$ a homogenous point~$\Epoint{\vectob{p}} \setin \reals^{n+1}$.
%
\begin{equation}
    \label{eq:homogenous}
    \Ematrix{\matmor{R}}{\vectmor{t}}
    \Epoint{\vectob{p}}
    =
    \Epoint{\matmor{R}\vectob{p} + \vectmor{t}}.
\end{equation}
%

If we apply two rototranslations, first~$\monela=\tupp{\matmor{R}_{\monela}, \vectmor{t}_{\monela}}$ and then~$\monelb=\tupp{\matmor{R}_{\monelb}, \vectmor{t}_{\monelb}}$, we find:
%
\begin{equation}
    \label{eq:rototranslation-series}
    \begin{aligned}
        \mgact(\tupp{\matmor{R}_{\monelb}, \vectmor{t}_{\monelb}}, \mgact(\tupp{\matmor{R}_{\monela}, \vectmor{t}_{\monela}}, \vectob{p})) & =
        \mgact(\tupp{\matmor{R}_{\monelb}, \vectmor{t}_{\monelb}}, \matmor{R}_{\monela}\vectob{p} + \vectmor{t}_{\monela}) \\
                                                                                                                                           & = \matmor{R}_{\monelb}\matmor{R}_{\monela}\vectob{p} + \matmor{R}_{\monelb}\vectmor{t}_{\monela} + \vectmor{t}_{\monelb}.
    \end{aligned}
\end{equation}

Alternatively, we can first compose the transformation~$\monela$ with~$\monelb$:
%
\begin{equation}
    \label{eq:rototranslation-series-2}
    \tupp{\matmor{R}_{\monela}, \vectmor{t}_{\monela}} \mtimesof{\mgSEn} \tupp{\matmor{R}_{\monelb}, \vectmor{t}_{\monelb}} = \tupp{\matmor{R}_{\monelb} \matmor{R}_{\monela}, \matmor{R}_{\monelb} \vectmor{t}_{\monela} + \vectmor{t}_{\monelb}},
\end{equation}
%
and then apply this composite transformation to the point~$\vectob{p}$:
\begin{equation}
    \label{eq:rototranslation-series-3}
    \mgact(\tupp{\matmor{R}_{\monelb} \matmor{R}_{\monela}, \matmor{R}_{\monelb} \vectmor{t}_{\monela} + \vectmor{t}_{\monelb}}, \vectob{p})
    = \matmor{R}_{\monelb}\matmor{R}_{\monela}\vectob{p} + \matmor{R}_{\monelb}\vectmor{t}_{\monela} + \vectmor{t}_{\monelb}.
\end{equation}
%
Thus, we have proved this property
%
\begin{equation}
    \label{eq:rototranslation-action}
    \mgact(\monelb, \mgact(\monela, \vectob{p})) = \mgact(\monela \mtimes \monelb, \vectob{p}),
\end{equation}
which is graphically reported in \cref{fig:graphical-roto-action}.

\begin{figure}[h]
    \includesag{10_graph_roto_action}
    \caption{Graphical representation of roto-translation action.}
    \label{fig:graphical-roto-action}
\end{figure}
%
The notion of \SY{semigroup action} generalizes this property.





